\section{滞后回归与 ARDL：把回归变成动态模型}
\label{sec:06-Random-Processes/06-ARMA-and-ARIMA/05}

定价团队问了一个听上去很简单的问题：把某个品类的价格永久下调一档，销量会涨多少？

你手上有周度的价格和销量，跑一个回归 \(Y_t=\beta_0+\beta_1X_t+\varepsilon_t\)，拿到一个系数，报上去。一个季度之后复盘，实际涨幅比你报的数高出一大截。

回归没算错，它回答的是另一个问题。同期系数说的是``价格变动那一周，销量变动多少''；定价团队想知道的是``尘埃落定之后，销量停在什么水平''。中间隔着几周的调整过程：渠道要补货，用户要发现，销量自己还有惯性——上周卖得多，这周的基数就高。静态回归里没有任何位置容纳这些。

\subsection{把时间放进回归元}

补法很直接，把滞后项加进右边：

\[
Y_t=c+\phi_1Y_{t-1}+\beta_0X_t+\beta_1X_{t-1}+\varepsilon_t.
\]

这是自回归分布滞后模型（ARDL）最小的样子：因变量带自回归项，外生变量带分布滞后项。一般形式是

\[
Y_t=c+\sum_{i=1}^{p}\phi_iY_{t-i}+\sum_{j=0}^{q}\beta_jX_{t-j}+\varepsilon_t,
\]

记作 ARDL(\(p,q\))。全部 \(\beta_j\) 为零时它退回 AR(\(p\))；全部 \(\phi_i\) 为零时它退回纯分布滞后模型。若还想让误差项本身带 MA 结构，那是 ARMAX 或者说动态回归，与 ARDL 不是一回事——ARDL 的误差被假设为白噪声，这条假设后面要单独付账。

\subsection{同一个系数回答不了两个问题}

现在把上面那个定价问题算出来。设 \(\phi_1=0.6\)，\(\beta_0=2\)，\(\beta_1=1\)。让 \(X\) 在 \(t=0\) 永久上升一个单位，之前一直为零，看 \(Y\) 怎么走。

第 \(0\) 期只有 \(\beta_0\) 起作用，\(Y\) 跳 \(2\)。第 \(1\) 期 \(\beta_1\) 进来，同时上一期的 \(2\) 被 \(\phi_1\) 传递过来，\(Y=0.6\times2+2+1=4.2\)。此后 \(X\) 不再变化，每一期都是把上期水平乘 \(0.6\) 再加上 \(3\)，序列单调爬向不动点 \(3/(1-0.6)=7.5\)。

\begin{center}
\begin{tikzpicture}[x=0.62cm,y=0.32cm,font=\small,>=stealth]
  \draw[->,black!55] (-2.8,0) -- (13.4,0) node[right,font=\footnotesize] {\(t\)};
  \draw[->,black!55] (-2.8,0) -- (-2.8,9.4);
  \draw[black!30,dashed] (-2.8,7.5) -- (12.6,7.5);
  \draw[black!30,dashed] (-2.8,2) -- (12.6,2);
  \draw[black!30,dashed] (-2.8,4.75) -- (3.0,4.75);
  \draw[black!45,dashed] (0,-0.8) -- (0,8.8);
  \node[above,font=\footnotesize,black!65,right] at (0.3,8.5) {\(X\) 在此永久上升一个单位};
  \draw[blue!70,thick] (-2.8,0) -- (0,0) -- (0,2) -- (1,4.2) -- (2,5.52) -- (3,6.31) -- (4,6.79)
    -- (5,7.07) -- (6,7.24) -- (7,7.35) -- (8,7.41) -- (9,7.45) -- (10,7.47) -- (11,7.48) -- (12,7.49);
  \foreach \p in {(0,2),(1,4.2),(2,5.52),(3,6.31),(4,6.79),(5,7.07),(6,7.24),(7,7.35)}
    \fill[blue!75] \p circle (1.6pt);
  \node[right,font=\footnotesize,black!70] at (12.7,7.5) {长期 \(7.5\)};
  \node[right,font=\footnotesize,black!70] at (12.7,2.0) {当期 \(2\)};
  \node[right,font=\footnotesize,black!55] at (3.1,4.75) {走完一半只用一期多一点};
\end{tikzpicture}

\emph{一次永久性的输入变化，激起一整条调整路径。当期效应 \(\beta_0=2\) 只是这条路径的起点，终点是长期乘数 \(7.5\)，两者差了将近四倍。静态回归的系数落在起点上，业务方问的却是终点。}
\end{center}

这条路径的两个端点各有名字。当期效应就是 \(\beta_0\)。长期乘数是

\[
\frac{\sum_{j=0}^{q}\beta_j}{1-\sum_{i=1}^{p}\phi_i},
\]

分子是外生变量所有滞后的直接贡献之和，分母是自回归反馈的放大倍数。分母越接近零——也就是 \(\sum\phi_i\) 越接近 \(1\)——放大得越厉害，调整也越慢。这里再次遇到第一节那道墙：\(\sum\phi_i=1\) 时分母为零，长期乘数不存在，因为系统根本不收敛到任何水平。所以``长期效应是多少''这个问题只在自回归部分稳定时才有答案。

调整速度也能读出来。偏离长期水平的部分每期乘 \(\phi_1\)，走完一半需要 \(\log0.5/\log\phi_1\) 期，这里约 \(1.4\) 期。报告一个长期乘数而不报告调整速度，等于告诉对方终点却不说要走多久。

\subsection{误差修正只是同一个模型换个写法}

把 ARDL(1,1) 两边同减 \(Y_{t-1}\) 并重新配项：

\[
\Delta Y_t=\beta_0\Delta X_t
+(\phi_1-1)\left(Y_{t-1}-\frac{c}{1-\phi_1}-\frac{\beta_0+\beta_1}{1-\phi_1}X_{t-1}\right)+\varepsilon_t.
\]

这一步是纯代数，没有引入任何新假设。括号里的东西正是上一节算出的长期关系

\[
Y=\frac{c}{1-\phi_1}+\frac{\beta_0+\beta_1}{1-\phi_1}X
\]

的偏离量，而 \(\phi_1-1\) 是它每期被修正掉的比例。写成这个形式之后，模型读起来是一句话：本期的变化由两部分组成，一部分是对当期输入变化的直接反应，另一部分是对上期偏离均衡的回拉。所以人们叫它误差修正模型（ECM）。

修正系数的取值范围决定行为。\(0<\phi_1<1\) 时 \(\phi_1-1\) 落在 \((-1,0)\)，偏离被单调地拉回；\(-1<\phi_1<0\) 时 \(\phi_1-1\) 落在 \((-2,-1)\)，每次修正都会冲过头，围绕均衡振荡收敛；\(\phi_1\ge1\) 或 \(\phi_1\le-1\) 时没有回拉，``长期均衡''这个词失去意义。估出来的修正系数为正，或者绝对值大于 \(2\)，都是设定出了问题的信号，不是有趣的发现。

\subsection{协整：代数恒等式什么时候才是统计关系}

ECM 的重写永远成立，它的\emph{解释}却需要额外条件，这是滞后回归里最容易跳过的一步。

假设 \(Y_t\) 与 \(X_t\) 各自都有单位根。括号里的线性组合 \(Y_{t-1}-a-bX_{t-1}\) 一般来说仍然是单位根过程，它不会回到零附近，也就没有``偏离''可言。此时你写下的 ECM 在代数上无懈可击，在统计上是空的：回归得到的系数不收敛到任何有意义的量，\(t\) 统计量随样本量发散，拟合优度看起来漂亮得可疑。这就是上一节说过的伪回归，只不过换了个外衣。

只有当这个线性组合恰好平稳时——两条序列各自漫游，但它们之间的距离被拉住——才叫协整，ECM 的解读才成立。协整对应的是真实的共同驱动：短期利率与长期利率、同一商品在两个市场的价格、消费与收入。检验协整有 Engle--Granger 的两步法（先回归再检验残差是否平稳）和 Johansen 的系统方法（可以同时处理多个协整关系）。

判断次序值得固定下来：先看每条序列的单整阶数，再看它们的线性组合是否平稳，最后才谈长期均衡。反过来做——先跑回归，看见显著就宣布找到了长期关系——是这一类分析中最常见的事故。

\subsection{OLS 什么时候还能用}

滞后因变量做回归元，把 OLS 的适用条件收紧了。

若 \(\varepsilon_t\) 确实是白噪声且与所有回归元无关，OLS 一致且渐近正态，有限样本略有偏（因为 \(Y_{t-1}\) 与更早的 \(\varepsilon\) 相关），样本一长就可以忽略。麻烦出在误差有序列相关的时候：\(Y_{t-1}\) 里含着 \(\varepsilon_{t-1}\)，而 \(\varepsilon_{t-1}\) 与 \(\varepsilon_t\) 相关，于是回归元与误差项相关，OLS 不但有偏，而且不一致——样本再多也不会收敛到真值。加数据在这里救不了。

诊断上有个具体的坑：Durbin--Watson 统计量在存在滞后因变量时会系统性地偏向``无自相关''，用它做检验等于自己骗自己。该用的是 Breusch--Godfrey 检验，或者直接看残差的自相关图。发现序列相关之后，两条出路：把误差结构显式建进去（ARMAX、或者干脆加更多 \(Y\) 的滞后项把相关吸收掉），或者改用工具变量与广义矩估计。

另有一层假设常被默认为真，就是 \(X_t\) 的外生性。若 \(X\) 与 \(Y\) 互相影响——价格影响销量，销量又反过来影响下一期的定价——那 \(X_t\) 与 \(\varepsilon_t\) 相关，系数估计不再是因果效应。此时该用的是向量自回归那套对称的框架，把两条序列放在同等地位上。

\subsection{加了滞后项，仍然不是因果}

最后一条边界必须讲清楚，因为 ARDL 的形式太容易引人误读。

方程右边站着 \(X_{t-1}\) 而不是 \(X_{t+1}\)，这个时间上的先后关系让人产生一种错觉：既然原因在前，那估出来的系数就是因果效应。Granger 意义上的``因果''也强化了这个错觉，但那个词的定义只是``加入 \(X\) 的历史能不能显著改进对 \(Y\) 的预测''。预测得更准与干预会奏效是两回事。

反例不难构造。两条序列同受一个未观测的共同驱动影响，而 \(X\) 对这个驱动的反应快一点，\(Y\) 慢一点——比如宏观景气同时推动广告投放与销量，投放调整得快，销量反应得慢。那么 \(X\) 的历史确实能预测 \(Y\)，Granger 检验会显著，可停掉投放并不会让销量按系数所说的幅度下降。这正是\hyperref[sec:13-Machine-Learning/12-Causal-Inference/03]{``混杂、碰撞点与后门调整''}里那个共同原因结构，只是把它摊在了时间轴上。时间顺序能排除``果在因前''这一种错误，排除不了混杂。

于是滞后回归的合适定位是：一个把惯性、延迟与调整过程写清楚的预测与描述工具。它给出的短期效应、长期乘数、调整半衰期都是关于\emph{这套数据生成机制}的量化描述，只有在设计（随机化、自然实验）或者可信的识别假设支撑下，它们才能被读作干预的后果。

本单元到此走完一条线：从序列对自己回归，到冲击的有限记忆，到用差分把非平稳拉回框架内，再到把外部变量接进来。全程只用了一种数学结构——线性滞后算子的多项式，以及它的根落在单位圆的哪一侧。这个结构简单到可以完全掌握，也因此它的失效方式一条条都能说清楚：非线性、状态切换、条件异方差、参数漂移、混杂。知道一个模型在哪里失效，比知道它在哪里成立更接近可用的判断。
